% arquivo_principal.tex

\documentclass[12pt, a4paper]{article}
\usepackage [brazil]{babel} 
\usepackage[utf8]{inputenc} 
\usepackage[left=3cm,top=2cm,right=1.5cm,bottom=2cm]{geometry}
\usepackage{graphicx}
\usepackage{setspace}
\usepackage{mathrsfs}
\usepackage{amsmath}
\usepackage{float}
\usepackage{amssymb}
\usepackage{mathtools}
\usepackage{fontenc}
\usepackage{multirow} 
\usepackage[breaklinks=true]{hyperref}
\usepackage[table,xcdraw]{xcolor}
\usepackage{indentfirst}
\usepackage{listings} 
\usepackage[labelfont=bf]{caption} 
\usepackage{tikz}
\usetikzlibrary{shapes, arrows, positioning}

%%%%%%%%%%%%%%%%%%%%%%%%%
%% Início do Código  para Python
%%%%%%%%%%%%%%%%%%%%%%%%
\lstset{ %
	backgroundcolor=\color{gray!1},   
	basicstyle=\ttfamily\small,		 
	breaklines=true,                 
	captionpos=b,                    
	commentstyle=\color{green!60!black},    
	escapeinside={\%}{)},          
	keywordstyle=\color{magenta},       
	stringstyle=\color{blue!50!red},     
	showstringspaces=false,			 
	language=python,                 
	numbers=left,
	numberstyle=\footnotesize\color{gray},
	numbersep=10pt,
	tabsize=4,
	frame=l, 
}
%%%%%%%%%%%%%%%%%%%%%%%%%
%% Fim do Código  para Python
%%%%%%%%%%%%%%%%%%%%%%%%

\begin{document}
	
	\begin{titlepage}
		
		\begin{figure}[H]
			\centering
			\includegraphics[scale=0.4]{Figuras/LogoUesc.png}
			\label{logouesc}
		\end{figure}
		\begin{center} 
			\textbf{UNIVERSIDADE ESTADUAL DE SANTA CRUZ - UESC\\
				PRÓ-REITORIA DE PESQUISA E PÓS-GRADUAÇÃO \\
				PROGRAMA INSTITUCIONAL DE INICIAÇÃO CIENTÍFICA\\
				Rod. Jorge Amado, Km 16 CEP 45650-900 Ilhéus-Bahia}\\\vspace{3cm}
			
			
			\textbf{RELATÓRIO FINAL - FAPESB}\\\vspace{1cm}
			\textbf{\color{blue}Controle PID em Sistema de Aquecimento em Planta Química}\\\vspace{3cm}
			\hspace{100cm}
			
			\hspace{6cm}
			\begin{minipage}[t]{.6\textwidth}
				\small Relatório Final de Iniciação Científica ICB que foi realizado pelo discente \textbf{202210950} Gabriel Marques de Andrade à Universidade Estadual de Santa Cruz - UESC.\\ 
				
				\textbf{Orientador}: Prof. Dr. Elilton Rodrigues Edwards.
			\end{minipage}\\\vspace{5cm}
					
			Ilhéus, \today{}.
			
		\end{center}

\end{titlepage}

%%%%%%%%%%%%%%%%%%%%%%%%%%%%%%%%%%%%%%%%%%%%%%
% ESPAÇAMENTO DO TEXTO 
%%%%%%%%%%%%%%%%%%%%%%%%%%%%%%%%%%%%%%%%%%%%%%%
\doublespacing 

%%%%%%%%%%%%%%%%%%%%%%%%%%%%%%%%%%%%%%%%%%%%%%
% LISTA DE FIGURAS, TABELAS, ABREVIATURAS, ETC 
%%%%%%%%%%%%%%%%%%%%%%%%%%%%%%%%%%%%%%%%%%%%%%%

{\color{blue}\tableofcontents }
\pagebreak

%%%%%%%%%%%%%%%%%%%%%%%%%%%%%%%%%%%%%%%%%%%%%%%%%%%%%%%%%%%%%%%
%  RESUMO 
%%%%%%%%%%%%%%%%%%%%%%%%%%%%%%%%%%%%%%%%%%%%%%%%%%%%%%%%%%%%%%%
\section*{Resumo}

Este projeto propõe o desenvolvimento e a análise de estratégias de controle digital aplicadas a processos térmicos na indústria química, com foco em um Tanque de Aquecimento com Agitação Contínua (CSTH). A metodologia baseia-se na formulação de modelos fenomenológicos via balanços de massa e energia, seguida de sua discretização e implementação computacional utilizando a linguagem Python e bibliotecas científicas (SciPy e Matplotlib). O escopo abrange o projeto de controladores Proporcional-Integral-Derivativo (PID) digitais e a avaliação comparativa de métodos clássicos de sintonia (Ziegler-Nichols e Cohen-Coon) frente a abordagens avançadas baseadas em Controle por Modelo Interno (IMC). O desempenho do sistema em malha fechada será quantificado sob cenários de rastreamento de referência e rejeição a perturbações, visando maximizar a estabilidade e a robustez térmica do processo.

\textbf{Palavras-chave:} Controle de Processos; Sistema Térmico; Controlador PID Digital; Python; Sintonia IMC.

%%%%%%%%%%%%%%%%%%%%%%%%%%%%%%%%%%%%%%%%%%%%%%%%%%%%%%%%%%%%%%%
%                 OBJETIVO GERAL
%%%%%%%%%%%%%%%%%%%%%%%%%%%%%%%%%%%%%%%%%%%%%%%%%%%%%%%%%%%%%%%
\newpage
\section{Objetivo Geral}

Desenvolver a modelagem dinâmica de um sistema de aquecimento representativo de plantas químicas e implementar controladores PID digitais utilizando Python e suas bibliotecas científicas para controle de temperatura, garantindo estabilidade analítica e mitigação de distúrbios operacionais.

%%%%%%%%%%%%%%%%%%%%%%%%%%%%%%%%%%%%%%%%%%%%%%%%%%%%%%%%%%%%%%%
%                 OBJETIVOS ESPECÍFICOS
%%%%%%%%%%%%%%%%%%%%%%%%%%%%%%%%%%%%%%%%%%%%%%%%%%%%%%%%%%%%%%%

\section{Objetivos Específicos}

\begin{itemize}
    \item Formular o modelo matemático do processo térmico por meio de equações diferenciais ordinárias fundamentadas em balanços de massa e energia.
    \item Realizar a discretização do modelo contínuo para viabilizar a implementação estrita de algoritmos de controle em domínio digital.
    \item Desenvolver o ambiente de simulação em Python, utilizando a biblioteca SciPy para a integração numérica da dinâmica da planta.
    \item Projetar e programar controladores nas topologias P, PI e PID no laço de realimentação.
    \item Aplicar e comparar métodos de sintonia clássicos (Ziegler-Nichols, Cohen-Coon) e o método analítico avançado de Controle por Modelo Interno (IMC).
    \item Analisar o desempenho, a estabilidade e a robustez do sistema de controle sob condições de mudança de setpoint, perturbações externas e variações paramétricas na carga térmica.
\end{itemize}

%%%%%%%%%%%%%%%%%%%%%%%%%%%%%%%%%%%%%%%%%%%%%%%%%%%%%%%%%%%%%%%
%                 INTRODUÇÃO
%%%%%%%%%%%%%%%%%%%%%%%%%%%%%%%%%%%%%%%%%%%%%%%%%%%%%%%%%%%%%%%
\newpage
\section{Introdução}

A regulação de variáveis termodinâmicas em estado estacionário ou transiente é um requisito crítico na engenharia de processos químicos, exercendo impacto direto sobre a eficiência energética, segurança operacional e cinética reacional das plantas industriais \cite{sen2004}. Dentre os equipamentos de transferência de calor, o Tanque de Aquecimento com Agitação Contínua (CSTH) se consolida como um paradigma complexo de controle devido à sua dinâmica não-linear inerente e ao forte acoplamento existente entre a vazão dos fluidos e a transferência de entalpia \cite{dhanasekar2024}.

O controlador Proporcional-Integral-Derivativo (PID) domina a automação regulatória industrial devido à sua simplicidade, robustez e eficácia comprovada em um amplo espectro de processos, incluindo controle de temperatura, pressão e fluxo \cite{stanisic2025}. O algoritmo de controle PID baseia-se na manipulação do erro, definido como a diferença entre a variável de processo e o \textit{setpoint}, e é composto pela soma de três termos: um termo proporcional ao erro, um termo proporcional à integral do erro e um termo proporcional à derivada do erro \cite{khare2010}. A representação no domínio do tempo da lei de controle PID ideal é expressa pela Equação \ref{eq:pid_ideal}:

\begin{equation}u(t) = K_p \left( e(t) + \frac{1}{T_i} \int_{0}^{t} e(\tau) d\tau + T_d \frac{de(t)}{dt} \right)\label{eq:pid_ideal}\end{equation}\noindent onde $u(t)$ é o sinal de controle, $e(t)$ é o erro, $K_p$ é o ganho proporcional, $T_i$ é o tempo integral e $T_d$ é o tempo derivativo \cite{khare2010}.

A topologia em malha fechada para o controle térmico exige a aquisição contínua da temperatura e a atuação correspondente sobre a variável manipulada, tipicamente a vazão do fluido de aquecimento (Figura \ref{fig:diagrama_blocos}).

\begin{figure}[H]
    \centering
    \begin{tikzpicture}[auto, node distance=1.2cm, >=latex']
        % Definição de estilos locais
        \tikzset{
            block/.style = {draw, rectangle, minimum height=3em, minimum width=4em, align=center},
            sum/.style = {draw, circle, node distance=1.2cm},
            input/.style = {coordinate},
            output/.style = {coordinate}
        }
        
        % Inserção dos Nós
        \node [input, name=input] {};
        \node [sum, right=of input] (sum1) {$\Sigma$};
        \node [block, right=of sum1] (pid) {Controlador\\PID};
        \node [block, right=of pid, node distance=1.8cm] (atuador) {Atuador/\\Válvula};
        \node [sum, right=of atuador] (sum2) {$\Sigma$};
        \node [block, right=of sum2] (planta) {Processo\\(CSTH)};
        \node [output, right=of planta, node distance=1.8cm] (output) {};
        \node [input, above=of sum2, node distance=1.5cm] (disturb) {};
        \node [block, below=of atuador, node distance=1.5cm] (sensor) {Sensor/\\Transmissor};

        % Conexões (Setas)
        \draw [->] (input) -- node {$T_{ref}(t)$} (sum1);
        \draw [->] (sum1) -- node {$e(t)$} (pid);
        \draw [->] (pid) -- node {$u(t)$} (atuador);
        \draw [->] (atuador) -- (sum2);
        \draw [->] (sum2) -- (planta);
        \draw [->] (planta) -- node [name=y] {$T(t)$} (output);
        \draw [->] (y) |- (sensor);
        \draw [->] (sensor) -| node[pos=0.95, right] {$-$} (sum1);
        \draw [->] (disturb) -- node[pos=0.2, right] {Perturbação $d(t)$} (sum2);
    \end{tikzpicture}
    \caption{Diagrama de blocos conceitual de um sistema de controle de temperatura em malha fechada com rejeição a perturbações.}
    \label{fig:diagrama_blocos}
\end{figure}

Historicamente implementados sob lógicas analógicas, as atuais arquiteturas industriais exigem a conversão destes sistemas para topologias de controle digital. A implementação no domínio discreto permite maior flexibilidade algorítmica, tratamento de restrições operacionais e a integração com estratégias avançadas de otimização dinâmica \cite{hedengren2014}. A aplicação de controladores PID digitais, aliada a simulações computacionais baseadas em modelos rigorosos de equações diferenciais e algébricas, configura-se como um vetor necessário para a validação da estabilidade estrutural e análise de robustez térmica frente a perturbações não mensuráveis \cite{hedengren2014}.

Este projeto fundamenta-se na necessidade técnica de analisar rigorosamente a dinâmica do processo térmico para estabelecer métricas comparativas entre algoritmos de sintonia analítica convencionais e técnicas avançadas baseadas em modelo (Controle por Modelo Interno - IMC), suprimindo empirismos na determinação de ganhos operacionais e garantindo o rastreamento de referência com eliminação de erro em regime permanente \cite{stanisic2025}.

%%%%%%%%%%%%%%%%%%%%%%%%%%%%%%%%%%%%%%%%%%%%%%%%%%%%%%%%%%%%%%%
%                 REVISÃO BIBLIOGRÁFICA
%%%%%%%%%%%%%%%%%%%%%%%%%%%%%%%%%%%%%%%%%%%%%%%%%%%%%%%%%%%%%%%
\newpage
\section{Revisão Bibliográfica}

A precisão do projeto de uma malha de controle exige modelagem fenomenológica rigorosa do sistema em estudo. Sistemas térmicos reais demonstram comportamento com atrasos de transporte significativos (tempo morto) devido ao deslocamento do fluido e inércia térmica dos metais, o que compromete severamente a margem de fase do sistema em malha fechada \cite{sen2004}. Conforme delineado por \cite{dhanasekar2024}, a dinâmica de um CSTH requer a resolução estrita de equações de balanço energético. Para viabilizar a derivação algébrica de ganhos de controle, a dinâmica complexa é frequentemente reduzida a um modelo de Primeira Ordem com Atraso de Transporte (FOPTD - \textit{First-Order Plus Dead-Time}), cuja função de transferência empírica é expressa pela Equação \ref{eq:foptd}:

\begin{equation}
    G(s) = \frac{K}{\tau s + 1} e^{-\theta s}
    \label{eq:foptd}
\end{equation}

\noindent onde $K$ é o ganho estático do processo, $\tau$ é a constante de tempo principal e $\theta$ é o tempo morto (atraso de transporte) \cite{skoge2003}.

Em nível industrial, o controlador PID domina a automação regulatória devido à sua arquitetura reativa, integrativa e antecipatória simples \cite{stanisic2025}. A parametrização dos ganhos ($K_p$, $T_i$, $T_d$) depende da correta identificação dos parâmetros do modelo FOPTD ($K, \tau, \theta$). Historicamente, as formulações de Ziegler-Nichols (Z-N) e Cohen-Coon baseadas na curva de reação ao degrau em malha aberta têm sido o padrão da indústria. A Tabela \ref{tab:sintonia} consolida as regras de sintonia clássicas \cite{odwyer2007}. 

Entretanto, literatura recente comprova que, diante de processos com predominância de tempo morto ($\theta \approx \tau$), sintoniais clássicas introduzem elevado sobressinal (\textit{overshoot}) e baixa margem de robustez \cite{khare2010}. Para solucionar tais limitações, o Controle por Modelo Interno (IMC), especificamente através das regras analíticas SIMC (\textit{Skogestad IMC}), fornece uma estrutura de parametrização baseada na inversão do modelo. O método SIMC substitui o empirismo por um parâmetro de projeto ($\tau_c$), que define a constante de tempo desejada para a malha fechada, estabelecendo um compromisso exato entre velocidade de resposta e robustez \cite{skoge2003, padhee2011}. As equações de sintonia analítica IMC para a topologia PID também estão sumarizadas na Tabela \ref{tab:sintonia}.

\begin{table}[H]
    \centering
    \caption{Correlações analíticas para sintonia de controladores PID baseadas no modelo FOPTD.}
    \label{tab:sintonia}
    \renewcommand{\arraystretch}{1.8} 
    \resizebox{\textwidth}{!}{%
    \begin{tabular}{|c|c|c|c|}
        \hline
        \textbf{Método} & \textbf{Ganho Proporcional ($K_p$)} & \textbf{Tempo Integral ($T_i$)} & \textbf{Tempo Derivativo ($T_d$)} \\ \hline
        Ziegler-Nichols & $\dfrac{1.2 \tau}{K \theta}$ & $2 \theta$ & $0.5 \theta$ \\ \hline
        Cohen-Coon & $\dfrac{\tau}{K \theta} \left( \dfrac{4}{3} + \dfrac{\theta}{4\tau} \right)$ & $\theta \left( \dfrac{32 + 6\theta/\tau}{13 + 8\theta/\tau} \right)$ & $\theta \left( \dfrac{4}{11 + 2\theta/\tau} \right)$ \\ \hline
        IMC (SIMC) & $\dfrac{\tau}{K (\tau_c + \theta)}$ & $\min\{\tau, 4(\tau_c + \theta)\}$ & $\dfrac{\tau_2}{1 + \tau_2/\tau_1} \approx 0$ \\ \hline
    \end{tabular}%
    }
\end{table}

Hedengren et al. \cite{hedengren2014} reiteram que a discretização de modelos contínuos, integrada a resolvedores numéricos de equações diferenciais (DAEs), é uma transição metodológica indispensável na implementação, simulação e avaliação destas regras analíticas em ambientes de controle digital e otimização dinâmica.

%%%%%%%%%%%%%%%%%%%%%%%%%%%%%%%%%%%%%%%%%%%%%%%%%%%%%%%%%%%%%%%
%                 METODOLOGIA
%%%%%%%%%%%%%%%%%%%%%%%%%%%%%%%%%%%%%%%%%%%%%%%%%%%%%%%%%%%%%%%
\section{Metodologia}

Para garantir rigor analítico e reprodutibilidade de dados, a arquitetura deste projeto é dividida em cinco fases operacionais, executadas sistematicamente através do ambiente Python:

\begin{enumerate}
    \item \textbf{Modelagem Fenomenológica:} Estruturação das equações diferenciais ordinárias (EDOs) do tanque de aquecimento via balanços de massa global e energia, obtendo-se as taxas de variação térmica transiente do sistema de acordo com a vazão volumétrica manipulada.
    \item \textbf{Integração e Discretização Numérica:} Codificação do modelo analítico contínuo para o domínio digital, empregando o pacote \texttt{scipy.integrate} para garantir a precisão no cálculo das respostas temporais, assumindo tempos de amostragem condizentes à instrumentação simulada.
    \item \textbf{Projeto e Implementação da Malha de Controle:} Programação da equação em diferenças do controlador PID digital na malha de realimentação negativa, aplicando termos P, PI e PID conforme a ação do erro computado entre o \textit{setpoint} e a saída iterativa do integrador numérico.
    \item \textbf{Derivação de Sintonia Analítica e Heurística:} Execução do ensaio degrau em malha aberta no modelo computacional para identificação da planta FOPTD. Aplicação dos equacionamentos de Ziegler-Nichols, Cohen-Coon e o desenvolvimento analítico das regras de sintonia via IMC baseadas na inversão do modelo.
    \item \textbf{Avaliação de Resposta e Desempenho:} Simulação de cenários de resposta sob mudança degrau na referência e introdução arbitrária de ruído nas variáveis de perturbação (ex. fluxo e temperatura da corrente de entrada). Utilização da biblioteca \texttt{matplotlib} para geração das curvas operacionais e quantificação matemática da performance métrica por Índices de Erro (IAE, ISE, ITAE), além da documentação do \textit{overshoot} e tempo de acomodação.
\end{enumerate}

%%%%%%%%%%%%%%%%%%%%%%%%%%%%%%%%%%%%%%%%%%%%%%%%%%%%%%%%%%%%%%%
%                  RESULTADOS E DISCUSSÃO
%%%%%%%%%%%%%%%%%%%%%%%%%%%%%%%%%%%%%%%%%%%%%%%%%%%%%%%%%%%%%%%
\section{Resultados e Discussão}

(Os dados de integração numérica, comparação dos métodos de sintonia e plotagens das funções de transferência serão incorporados nesta seção nas fases subsequentes da execução do projeto).

%%%%%%%%%%%%%%%%%%%%%%%%%%%%%%%%%%%%%%%%%%%%%%%%%%%%%%%%%%%%%%%
%                  CONCLUSÃO
%%%%%%%%%%%%%%%%%%%%%%%%%%%%%%%%%%%%%%%%%%%%%%%%%%%%%%%%%%%%%%%
\section{Conclusão}

(A consolidação final referente ao comportamento dinâmico do CSTH sob os controladores parametrizados será detalhada após a simulação paramétrica completa).

%%%%%%%%%%%%%%%%%%%%%%%%%%%%%%%%%%%%%%%%%%%%%%%%%%%%%%%%%%%%%%%
%                  BIBLIOGRAFIA
%%%%%%%%%%%%%%%%%%%%%%%%%%%%%%%%%%%%%%%%%%%%%%%%%%%%%%%%%%%%%%%
\newpage
\bibliographystyle{abntex2-alf}
\bibliography{bibliografia} 

\end{document}